\documentclass[12pt,a4paper,openany]{report}
\usepackage[a4paper,margin=1in]{geometry}
\usepackage[T1]{fontenc}
\usepackage[utf8]{inputenc}
\usepackage{mathptmx}
\usepackage{setspace}
\usepackage{graphicx}
\usepackage{float}
\usepackage{booktabs}
\usepackage{array}
\usepackage{tabularx}
\usepackage{longtable}
\usepackage{xcolor}
\usepackage{hyperref}
\usepackage{caption}
\usepackage{tikz}
\usetikzlibrary{arrows.meta,positioning,shapes.geometric}
\usepackage{listings}
\usepackage{fancyhdr}
\usepackage{titlesec}
\usepackage{enumitem}
\usepackage{url}
\usepackage{placeins}

\hypersetup{colorlinks=true, linkcolor=black, urlcolor=blue, citecolor=black}
\onehalfspacing
\raggedbottom
\setlength{\parindent}{0pt}
\setlength{\parskip}{6pt}
\renewcommand{\arraystretch}{1.15}
\captionsetup{font=bf, labelfont=bf}
\urlstyle{same}

\definecolor{codegray}{RGB}{246,248,250}
\definecolor{codeborder}{RGB}{210,215,220}
\lstdefinestyle{reportcode}{
  backgroundcolor=\color{codegray},
  basicstyle=\ttfamily\small,
  breaklines=true,
  frame=single,
  framerule=0.4pt,
  rulecolor=\color{codeborder},
  columns=fullflexible,
  keepspaces=true,
  showstringspaces=false,
  tabsize=2
}
\lstset{style=reportcode}

\pagestyle{fancy}
\fancyhf{}
\fancyhead[L]{23AIML049}
\fancyhead[R]{Agentic Requirement Analyzer -- Internship Report}
\fancyfoot[L]{CSPIT, CHARUSAT}
\fancyfoot[C]{\thepage}
\fancyfoot[R]{Department of AIML}
\renewcommand{\headrulewidth}{0.4pt}
\renewcommand{\footrulewidth}{0.4pt}
\setlength{\headheight}{15pt}

\titleformat{\chapter}[display]
  {\normalfont\bfseries\centering\Large}
  {CHAPTER \thechapter}
  {0.5em}
  {\Large}
\titlespacing*{\chapter}{0pt}{-5pt}{14pt}
\titleformat{\section}{\normalfont\bfseries\large}{\thesection}{1em}{}
\titleformat{\subsection}{\normalfont\bfseries}{\thesubsection}{1em}{}

\newcommand{\reporttitle}{AGENTIC REQUIREMENT ANALYZER}
\newcommand{\subtitle}{Autonomous Multi-Agent AI System for Software Requirement Analysis and User Story Generation}
\newcommand{\student}{HARI PATEL}
\newcommand{\rollno}{23AIML049}
\newcommand{\internalguide}{Prof. Hardik Jaiswal}
\newcommand{\externalguide}{Mr. Pallav Mamtora}
\newcommand{\externaltitle}{CEO}
\newcommand{\company}{Techmicra IT Solutions}
\newcommand{\hod}{Dr. Nirav Bhatt}


\begin{document}
\pagestyle{fancy}

% ===== original PDF pages 23-24 (report pages 14-15) =====
\setcounter{page}{14}
\setcounter{chapter}{5}\setcounter{section}{0}\setcounter{table}{0}\setcounter{figure}{0}
\chapter{CONCLUSION AND FUTURE SCOPE}

\section{Conclusion}
The objectives defined in Section 1.4 were achieved during the internship.
The modular six-stage analysis pipeline, implemented as a compiled
LangGraph \texttt{StateGraph}, together with multi-format document parsing,
actor/feature/constraint extraction, functional and non-functional
classification, user-story generation, acceptance-criteria generation,
FAISS-based ambiguity detection, rule-based risk detection, FastAPI
integration, React report rendering and SQLite persistence were fully
implemented, tested and demonstrated.

During implementation and review, I faced several technical challenges
that improved the quality of the project. I identified that the FAISS
\texttt{IndexFlatIP} vectors were not explicitly normalized, which meant
the scores could not safely be interpreted as cosine similarity, and that
the intended similarity threshold was described but not enforced before the
LLM confirmation step. In the frontend, acceptance-criteria fields already
contained the words ``Given'', ``When'' and ``Then'', while the renderer
added the same labels again. I fixed these issues, together with the
upload-endpoint weaknesses, and covered the fixes with automated tests.
The 0.40 similarity threshold was chosen empirically on three documents;
calibration on a labelled dataset remains future work.

Personally, the internship strengthened my ability to design agent
workflows, engineer prompts with strict JSON contracts, and implement
retry-once-then-fallback behaviour around unreliable model responses. I
gained practical experience with LangGraph orchestration, FastAPI endpoint
design, Sentence-Transformers embeddings, FAISS indexing, React integration
and persistent execution history. These skills connect directly with my
AIML coursework in natural language processing, machine learning
representations, software engineering and intelligent systems, while also
showing how model reasoning must be combined with validation, testing and
backend architecture to build a usable AI application.

% ----- original PDF page 24 -----
\clearpage
\subsection{Limitations}
\begin{itemize}[leftmargin=*]
\item The current application has no authentication or authorization layer,
  so reports and uploaded documents are not protected by user identity or
  role-based access control.
\item The complete analysis pipeline runs synchronously inside the HTTP
  request, which limits concurrency and can cause long-running requests or
  timeouts.
\item SQLite persistence and the frontend are designed for a single-user
  demonstration and do not provide tenant separation, collaborative review
  or multi-user report ownership.
\item Risk detection is rule-based rather than model-based, and the
  ambiguity threshold has not been validated on a labelled dataset.
\end{itemize}

\section{Future Scope}
Beyond the scope of this internship, the system could be extended with Jira
or Azure DevOps integration to push generated user stories directly into a
sprint board, UML and database schema generation from the same parsed
requirements, multi-language requirement support, and an effort-estimation
agent that assigns story points based on similar past requirements.
Production hardening should additionally apply secure file-upload controls,
API-security practices and documented AI risk-management procedures in
accordance with OWASP and NIST guidance [17,18,19].
\end{document}
